\subsection{攻读博士学位期间取得的学术成果}

期刊格式：
[序号] 作者. 论文名称[J]. 期刊名称，年度，卷（期）：起止页码. （检索情况）（对应论文章
节）

专利格式：

[序号] 专利申请者. 专利题名：专利国别，专利号[P]. 发布日期. （对应论文章节）

示例：

[1] Huang S C, Huang Y M, Shieh S M. Vibration and stability of a rotating shaft containing a transerse crack[J]. J Sound and Vibration, 1993, 162(3): 387-401. （SCI检索）（对应论文第四章）

[2] 高航，张立成，周士昌. 高压辊磨机液压系统及其动态特性[J]. 东北大学学报，2000，21（1）：38-40. （EI检索）（对应论文第五章）

[3] 刘加林. 多功能一次性压舌板：中国，92214985.2[P]. 1993-04-14. （对应论文第四章）


注：双盲评审版学位论文中须隐去所有作者（申请者）姓名，仅标注排序即可。

示例：

[1] 第一作者. Vibration and stability of a rotating shaft containing a transerse crack[J]. J Sound and Vibration, 1993, 162(3): 387-401. （SCI检索）（对应论文第四章）

[2] 第二作者. 高压辊磨机液压系统及其动态特性[J]. 东北大学学报，2000，21（1）：38-40. （EI检索）（对应论文第五章）

[3] 第二排序. 多功能一次性压舌板：中国，92214985.2[P]. 1993-04-14. （对应论文第四章）